

\chapter*{ABSTRACT}
\phantomsection
\addcontentsline{toc}{chapter}{ABSTRACT}


\bigskip

A Brain-Computer Interface (BCI) lets a person control a computer or device directly using brain signals which are recorded from the scalp using EEG electrodes. In this project, a low cost EEG based BCI typing system is designed using eye blinks and beta brain waves as control signals unlike complex systems that uses visual stimuli or machine learning, this system focuses on simplicity, affordability and ease of use. EEG signals are collected using electrodes and processed with basic signal processing techniques. Eye blinks create large changes in the EEG signal and are used for traversing the key in on-screen typing interface and focus state of brain can select the key. Waves generated from the brain sensed by electrodes are passed through analog filter and amplifiers to acquire adequate signal for computation of Fast Fourier Transform (FFT) on arduino microcontroller, where the obtained frequency information is used to differentiate the eye blink and focus state of brain which is the basis of using on-screen typing interface for our BCI system. A simple stimulus based method is able map these signals to typing and selection of commands  for interfacing. BCI can be helpful for people with movement difficulties and for situations where hands free control is preferred.


\vspace{1em}
\noindent\textit{\textbf{Keywords}: BCI, Electroencephalogram (EEG), Fast Fourier Transform (FFT), Signal Processing, Typing Interface}
